\documentclass[10pt]{article}

\usepackage[paperwidth=7in,paperheight=10in,margin=0.56in]{geometry}
\usepackage[no-math]{fontspec}
\usepackage{xcolor}
\usepackage{microtype}

\definecolor{Paleo}{HTML}{704214}
\newfontfamily\PaleoRegular[Path=../../fonts/,Renderer=HarfBuzz]{paleo-hebrew-regular.ttf}
\newfontfamily\PaleoEvyoni[Path=../../fonts/,Script=Hebrew,Renderer=HarfBuzz]{paleo-hebrew-evyoni.ttf}
\newfontfamily\PaleoAlphabet[Path=../../fonts/,Script=Hebrew,Renderer=HarfBuzz]{paleo-hebrew-alphabet.ttf}

\newcommand{\ASCIISample}[1]{{\PaleoRegular\fontsize{21}{25}\selectfont\color{Paleo}#1\par}}
\newcommand{\HebrewSample}[2]{{\textdir TRT#1\fontsize{19}{24}\selectfont\color{Paleo}#2\par}}
\newcommand{\SampleText}{אבגדהוזחטיכלמנסעפצקרשת}
\newcommand{\JacobText}{הקבצו בני יעקב ושמעו אל ישראל אביכם}

\pagestyle{empty}
\setlength{\parindent}{0pt}
\setlength{\parskip}{0.35em}

\begin{document}
{\Large Three new paleo-Hebrew hands}\par
{\small A proposal for the source typography of Genesis 49}\par

\section*{paleo-hebrew-regular.ttf --- \texttt{\string\hPaleoZ}}
\ASCIISample{ABGDHWZXJYKLMNS]PCQRVT}
\ASCIISample{HQBCW BNY Y]QB WVM]W AL YVRAL ABYKM}
This is the liveliest and most visibly hand-drawn face. Its irregular stroke and
open spacing suit reconstructed, orally inherited poetry, but it is a generic
modern evocation rather than an epigraphically dated hand. It could be useful
for a deliberately conjectural display, not as a replacement for the chapter's
dated Gezer, Tel Zayit, Mesha, or Siloam source intervals.

\section*{paleo-hebrew-evyoni.ttf --- \texttt{\string\hPaleoAA}}
\HebrewSample{\PaleoEvyoni}{\SampleText}
\HebrewSample{\PaleoEvyoni}{\JacobText}
This face is clean, compact, and readable in sustained Hebrew. It is the best
of the three for long reconstructed passages if legibility is allowed to outrank
the present inscription-based chronology. Its regularity, however, weakens the
chapter's useful claim that the changing hands carry historical argument.

\section*{paleo-hebrew-alphabet.ttf --- \texttt{\string\hPaleoAB}}
\HebrewSample{\PaleoAlphabet}{\SampleText}
\HebrewSample{\PaleoAlphabet}{\JacobText}
This face has the strongest monumental, incised character. That makes it a good
candidate for short headings, isolated tribal names, or a reconstructed line
presented as an emblem. Its uneven widths and emphatic silhouettes become noisy
across the long blessings, and its source is not a securely dated inscription.

\section*{Recommendation for 01-genesis/49.tex}
Keep the current choices for the chapter's historical source intervals. Gezer,
Tel Zayit, Mesha-style, and Siloam encode explicit chronological and editorial
claims which none of these generic faces can substantiate. If a distinct visual
register is later wanted for the most conjectural reconstructions, use
\texttt{\string\hPaleoZ} sparingly; if continuous-reading clarity becomes the
priority, test \texttt{\string\hPaleoAA}. Reserve \texttt{\string\hPaleoAB} for
short display matter. On the present evidence, no replacement in Genesis 49 is
better than the current historically named hands.
\end{document}
